\subsection*{Алгоритм Евклида}

\begin{lstlisting}
int gcd(int a, int b) {
    while (b != 0) {
        a %= b;
        swap(a, b);
    }
    return a;
}
\end{lstlisting}

\subsection*{Расширенный алгоритм  Евклида}

Расширенный алгоритм находит, помимо $g = gcd(a,\,b)$,  такие целые коэффициенты x и y, что $a \cdot x + b \cdot y = g$
Пусть мы посчитали нужные коэффициенты $x'\,\text{и}\,y'$, когда рекурсивно считали $gcd(b,\,a\,mod\,b)$. Иными словами, у нас есть решение $(x',\,y')$ для пары $(b,\,a\,mod\,b)$

$$b \cdot x' + (a\,mod\,b) \cdot y' = g$$

Чтобы получить решение для исходной пары, запишем выражение $(a\,mod\,b)$ по определению как $(a - \lfloor{\frac{a}{b}}\rfloor \cdot b)$ и подставим в приведенное выше равенство

$$b \cdot x' + (a - \lfloor{\frac{a}{b}}\rfloor \cdot b) \cdot y' = g$$

Теперь выполним перегруппировку слагаемых (сгруппируем по исходным a и b) и получим:

$$a \cdot y' + b \cdot (x' - \lfloor{\frac{a}{b}}\rfloor \cdot y') = g$$


$$\begin{cases}
x = y'\\
y = x' - \lfloor{\frac{a}{b}}\rfloor \cdot y'
\end{cases}$$

\begin{lstlisting}
int gcd(int a, int b, int &x, int &y) {
  if (b == 0) {
    x = 1;
    y = 0;
    return a;
  }
  int x1, y1;
  int d = gcd(b, a % b, x1, y1);
  x = y1;
  y = x1 - (a / b) * y1;
  return d;
}
\end{lstlisting}

Рассмотрим уравнение $ax + my = 1$. Возьмем обе части по модулю m
и $ax = 1 \Rightarrow x = a^{-1}$

То есть с помощью расширенного алгоритма Евклида ищем обратный
элемент в $\Z_m\,\text{за}\,O(log\,m)$.

Стоит отметить, что расширенный алгоритм Евклида не гарантирует, что $x < m$, поэтому нужно будет дополнительно сделать $x = x\; \%\; m$.

\pagebreak